\section{The AAT Approach}\label{sec:2}

Algebraic Architecture Theory (AAT) is a mathematical theory that
constructs software architecture as algebraic geometry. This section
briefly describes what is distinctive about the approach and the skeleton
of the construction. All definitions of the mathematical objects needed
for the statement and proof of SAGA are given in Section~\ref{sec:3}.

The construction proceeds in three stages. First, AAT axiomatizes
primitive architectural facts as \textbf{Atoms}. An Atom does not belong
to any particular programming language, framework, or serialization
format, so facts spanning multiple languages, services, and storage
representations can be treated on the same architecture object. To speak
about inconsistencies where the implementation types agree but the
semantic conventions differ, the architecture must be constructed at a
level of abstraction above the type system, and this language independence
is the prerequisite for the subject of this paper. Second, multiple
architecture conditions are organized as an \textbf{Atom-indexed
architectural equation system} $E$. From $E$ one constructs residuals,
the obstruction ideal, and the locus of points where the required
equations hold simultaneously; the place where the constraints are
satisfied becomes a geometric object, the analogue of the zero set of a
system of equations. Third, local contexts, covers, restrictions, and
overlaps are organized as an \textbf{AAT site}. The local-to-global
problem thereby becomes expressible in the language of sheaves and
cohomology: whether a family of local conditions glues globally is a
question about the existence of a global section, and the obstruction to
gluing is a class in \v{C}ech cohomology.

Under this construction, architectural failures are structured as
residuals, ideals, and cohomology classes, and diagnosis and repair can be
read as changes of the same geometric objects. A repair is an operation
that kills a class, and the success of a repair is judged by the vanishing
of the class. Furthermore, ringed geometry over an architecture opens
connections to the tools of algebraic geometry --- schemes, derived
intersections, deformations, monodromy --- but what SAGA uses in this
paper is \v{C}ech $H^1$ and descent; the deeper tools are discussed in the
research outlook of Section~\ref{sec:9}. The algebraic-geometric character
of SAGA lies on the side of the construction: the observable ring
$\ObsE$, the witness ideals and the obstruction ideal, the quotient
coefficient $\QE=\ObsE/\IOb$, the AAT site, and the descent
interpretation that reads the zero class as a global repair
(Sections~\ref{sec:3}--\ref{sec:5}; the attribution of what is proper to
SAGA is in \S\ref{sec:51}).

The geometry of AAT is a \textbf{relative geometry} with the vocabulary,
the equation system, the coverage, and the coefficients held fixed. Every
claim is relative to the selected input data, and the provenance from
inputs to conclusions can be traced. This relativity is a constraint, but
it is also the source of the discipline that matches claims with
evidence; the measurement of Section~\ref{sec:7} is its executable form.

SAGA concretizes the local-to-global capability of this approach by
comparing the semantic repair obstruction with the equation-generated
\v{C}ech obstruction. The approach of this section is summarized in one
sentence:

\begin{quote}
AAT constructs software architecture as a relative algebraic geometry
generated from primitive architectural facts and simultaneous
architectural equations, so that local consistency, global obstruction,
and repair become geometric objects.
\end{quote}
